\documentclass[12pt,letterpaper]{report}
\usepackage{graphicx}
\usepackage{scrextend}
\usepackage{vmargin}
\usepackage{graphicx}
\usepackage{multirow}
\usepackage[utf8]{inputenc}
\usepackage[spanish]{babel}
\usepackage{multicol}
\usepackage{enumerate}
\usepackage{float}
\usepackage{amsmath, amsthm, amssymb, amsfonts}
\usepackage[usenames]{color}
\parindent=0mm
\pagestyle{empty}
\definecolor{miorange}{rgb}{0.91, 0.43, 0.0}
\begin{document}
\setmargins{2.5cm}      
{1.5cm}                     
{2cm}  
{24cm}                    
{10pt}                          
{1cm}                          
{0pt}                             
{2cm}
\begin{flushright}
\today
\end{flushright}
\vspace{-1.75cm}
\section*{Tarea Clase 15}
\subsection*{Resuelva los siguientes ejercicios:}
\begin{enumerate}
    \item  Un dado está trucado, de forma que las probabilidades de obtener las distintas caras son proporcionales a los números de estas. Hallar:
    \begin{enumerate}
    \item La probabilidad de obtener el 6 en un lanzamiento
    \item La probabilidad de conseguir un número impar en un lanzamiento
    \end{enumerate}
    \item Se lanzan dos dados al aire y se anota la suma de los puntos obtenidos. Se pide:
    \begin{enumerate}
    \item La probabilidad de que salga el 7
    \item La probabilidad de que el número obtenido sea par
    \item La probabilidad de que el número obtenido sea múltiplo de tres 
    \end{enumerate}
    \item Se lanzan tres dados. Encontrar la probabilidad de que:
    \begin{enumerate}
    \item Salga 6 en todos
    \item Los puntos obtenidos sumen 7
    \end{enumerate}
    \item Busca la probabilidad de que al echar un dado al aire, salga:
    \begin{enumerate}
    \item Un número par
    \item Un múltiplo de tres
    \item Mayor que cuatro
    \end{enumerate}
    \item Se sacan dos bolas de una urna que se compone de una bola blanca, otra roja, otra verde y otra negra. Describir el espacio muestral cuando:
    \begin{enumerate}
    \item La primera bola se devuelve a la urna antes de sacar la segunda
    \item La primera bola no se devuelve
    \end{enumerate}
    \item Una urna contiene tres bolas rojas y siete blancas. Se extraen dos bolas al azar. Escribir el espacio muestral y hallar la probabilidad de:
    \begin{enumerate}
    \item Extraer las dos bolas con reemplazamiento
    \item Sin reemplazamiento
    \end{enumerate}
    \item En una clase hay 10 alumnas rubias, 20 morenas, cinco alumnos rubios y 10 morenos. Un día asisten 44 alumnos, encontrar la probabilidad de que el alumno que falta:
    \begin{enumerate}
    \item Sea hombre
    \item Sea mujer moreno
    \item Sea hombre o mujer
    \end{enumerate}
\end{enumerate}
\end{document}